% Analisis tujuh paper kandidat. Setiap paper dianalisis dengan delapan aspek yang sama
% supaya bisa dibandingkan. Isi diambil dari teks lengkap paper (versi arXiv dibaca, versi
% terbit yang dikutip); angka hasil disalin dari tabel paper masing-masing.

\section{Pendahuluan}\label{sec:pendahuluan}

Dokumen ini menganalisis tujuh paper kandidat untuk topik klasifikasi dokumen multimodal
yang memadukan citra halaman dengan teks hasil OCR. Paper dipilih dengan dua kriteria:
\begin{itemize}
  \item terbit tahun 2022 atau sesudahnya;
  \item terbit di konferensi atau jurnal yang disarankan.
\end{itemize}
Daftarnya pada Tabel~\ref{tab:kandidat}.

\begin{table}[H]
  \centering
  \caption{Paper kandidat yang dianalisis}
  \label{tab:kandidat}
  \fontsize{10}{12}\selectfont
  \begin{tabularx}{\textwidth}{|c|>{\raggedright\arraybackslash}X|l|c|}
    \hline
    \textbf{No} & \textbf{Paper} & \textbf{Venue} & \textbf{Tahun} \\ \hline
    1 & LayoutLMv3 \parencite{huang2022layoutlmv3} & ACM MM & 2022 \\ \hline
    2 & Donut \parencite{kim2022donut} & ECCV & 2022 \\ \hline
    3 & UDOP \parencite{tang2023udop} & CVPR & 2023 \\ \hline
    4 & DocFormerv2 \parencite{appalaraju2024docformerv2} & AAAI & 2024 \\ \hline
    5 & DocLLM \parencite{wang2024docllm} & ACL & 2024 \\ \hline
    6 & LiLT \parencite{wang2022lilt} & ACL & 2022 \\ \hline
    7 & Review fusion klasifikasi dokumen \parencite{mironczuk2026} &
        Information Fusion & 2026 \\ \hline
  \end{tabularx}
\end{table}

Setiap paper dianalisis dengan delapan aspek yang sama:
\begin{enumerate*}[label=(\arabic*)]
  \item latar belakang,
  \item dataset,
  \item teknologi OCR,
  \item pipeline,
  \item model,
  \item metode fusion,
  \item hasil, dan
  \item hambatan atau \textit{future work}.
\end{enumerate*}

Isinya diambil dari teks lengkap tiap paper. Yang dibaca adalah versi arXiv, tetapi yang
dikutip adalah versi terbitnya. Angka hasil disalin dari tabel paper masing-masing. Bila
sebuah aspek tidak disebutkan paper, hal itu ditulis \enquote{tidak disebutkan}, bukan
ditebak.

\section{Analisis per Paper}\label{sec:analisis}

\subsection{LayoutLMv3 (ACM MM 2022)}\label{sec:layoutlmv3}

\textbf{Latar belakang.} Model dokumen multimodal sebelumnya memakai tujuan pralatih yang
berbeda untuk teks dan citra:
\begin{itemize}
  \item teks memakai \textit{masked language modeling} (MLM);
  \item citra merekonstruksi piksel (DocFormer) atau meregresi fitur wilayah (SelfDoc).
\end{itemize}
Ketidakselarasan ini menyulitkan penyelarasan lintas-modalitas. Selain itu, cabang citra
bergantung pada CNN atau Faster R-CNN yang berat dan butuh anotasi wilayah
\parencite{huang2022layoutlmv3}.

\textbf{Dataset.}
\begin{itemize}
  \item \textbf{Pralatih:} IIT-CDIP, 11 juta halaman.
  \item \textbf{Klasifikasi:} RVL-CDIP (320.000/40.000/40.000 citra, 16 kelas).
  \item \textbf{Ekstraksi formulir dan struk:} FUNSD dan CORD.
  \item \textbf{Tanya-jawab dokumen:} DocVQA.
  \item \textbf{Analisis tata letak:} PubLayNet.
  \item \textbf{Lampiran:} EPHOIE (bahasa Mandarin).
\end{itemize}

\textbf{Teknologi OCR.}
\begin{itemize}
  \item Pralatih memakai \enquote{toolkit OCR siap pakai} yang tidak disebutkan namanya.
  \item Untuk RVL-CDIP dan DocVQA, teks dan kotaknya diekstrak dengan \textbf{Microsoft Read
  API} (layanan komersial).
  \item Untuk CORD dipakai anotasi OCR resmi dataset.
  \item Posisi tata letak dipakai di tingkat segmen: kata dalam satu segmen berbagi satu
  kotak.
\end{itemize}

\textbf{Pipeline.}
\begin{enumerate}
  \item Citra diubah ke $224\times224$, dipecah menjadi 196 \textit{patch} $16\times16$,
  lalu diproyeksikan linear.
  \item Teks OCR ditokenisasi BPE, maksimum 512 token, dan diberi posisi 1D serta posisi
  tata letak 2D (x, y, lebar, tinggi).
  \item Kedua urutan disambung dan dimasukkan ke satu Transformer.
  \item Untuk klasifikasi, keluaran token \texttt{[CLS]} masuk ke pengklasifikasi MLP.
\end{enumerate}

\textbf{Model.} Transformer \textit{encoder} yang diinisialisasi dari RoBERTa:
\begin{itemize}
  \item \textit{base}: 12 lapis, 133 juta parameter;
  \item \textit{large}: 24 lapis, 368 juta parameter.
\end{itemize}
Pralatihnya memakai tiga tujuan:
\begin{itemize}
  \item MLM untuk teks;
  \item \textit{masked image modeling} (MIM) untuk citra, dengan target token citra dari
  tokenizer DiT;
  \item \textit{word-patch alignment} (WPA): memprediksi apakah \textit{patch} citra dari
  sebuah kata ikut di-\textit{mask}.
\end{itemize}

\textbf{Metode fusion.} Fusion awal dalam satu aliran (\textit{single-stream early fusion}).
Token teks dan \textit{patch} citra disambung menjadi satu urutan, lalu saling memperhatikan
lewat \textit{self-attention} di setiap lapis. Penyelarasan antar-modalitas dilatih lewat
WPA.

\textbf{Hasil.}
\begin{itemize}
  \item RVL-CDIP: akurasi 95,44\% (\textit{base}) dan 95,93\% (\textit{large}).
  \item FUNSD: F1 90,29 dan 92,08. CORD: F1 96,56 dan 97,46.
  \item DocVQA: ANLS 78,76 dan 83,37.
  \item PubLayNet: mAP 95,1.
\end{itemize}

Ablasinya penting untuk topik kita. Pada pralatih 1 juta dokumen, model teks + tata letak
\textbf{tanpa citra} sudah mencapai 95,33\% di RVL-CDIP. Menambahkan \textit{patch} citra
tanpa MIM justru turun ke 95,00\%. Dengan MIM dan WPA, akurasinya baru naik ke 95,53\%.
Jadi sumbangan citra untuk klasifikasi hanya sekitar 0,2 poin, dan baru muncul setelah ada
tujuan pralatih khusus untuk citra.

\textbf{Hambatan dan \textit{future work}.} Penulis berencana:
\begin{itemize}
  \item memperbesar skala model dan data pralatih;
  \item mengeksplorasi \textit{few-shot} dan \textit{zero-shot learning}.
\end{itemize}
Hambatan yang tidak dibahas penulis:
\begin{itemize}
  \item model bergantung pada OCR komersial (Microsoft Read API);
  \item pralatihnya sangat mahal: \textit{batch} 2.048 selama 500.000 langkah.
\end{itemize}

\subsection{Donut (ECCV 2022)}\label{sec:donut}

\textbf{Latar belakang.} Model pemahaman dokumen menyerahkan pembacaan teks kepada mesin
OCR siap pakai. Ketergantungan ini membawa tiga masalah \parencite{kim2022donut}:
\begin{enumerate}
  \item biaya komputasi OCR yang tinggi;
  \item mesin OCR tidak fleksibel terhadap bahasa atau jenis dokumen baru;
  \item \textbf{galat OCR merambat} ke tahap berikutnya.
\end{enumerate}
Donut menghapus OCR dari pipeline.

\textbf{Dataset.}
\begin{itemize}
  \item \textbf{Pralatih:} IIT-CDIP (11 juta citra) ditambah dokumen sintetis SynthDoG,
  0,5 juta per bahasa (Inggris, Mandarin, Jepang, Korea).
  \item \textbf{Klasifikasi:} RVL-CDIP.
  \item \textbf{Ekstraksi informasi:} CORD, Ticket, serta dua data industri internal
  (kartu nama Jepang dan struk Korea).
  \item \textbf{Tanya-jawab:} DocVQA.
\end{itemize}

\textbf{Teknologi OCR.}
\begin{itemize}
  \item \textbf{Saat inferensi tidak ada OCR.}
  \item OCR hanya dipakai untuk membuat label teks pralatih IIT-CDIP, dengan CLOVA OCR API
  (komersial).
  \item Model pembanding berbasis OCR memakai MS OCR dan CLOVA OCR.
  \item Lampiran A.1 membandingkan \textbf{empat mesin OCR}: MS OCR, CLOVA OCR, EasyOCR, dan
  PaddleOCR versi \textit{mobile}.
\end{itemize}

\textbf{Pipeline.}
\begin{enumerate}
  \item Citra beresolusi $2560\times1920$ dimasukkan ke \textit{encoder} visual.
  \item \textit{Decoder} teks menghasilkan urutan token yang diubah menjadi JSON.
  \item Untuk klasifikasi, \textit{decoder} menuliskan \texttt{[START\_class][memo][END\_class]}.
\end{enumerate}

\textbf{Model.} 143 juta parameter:
\begin{itemize}
  \item \textit{encoder}: Swin Transformer-B;
  \item \textit{decoder}: empat lapis pertama BART multibahasa.
\end{itemize}
Pralatihnya berupa tugas \enquote{membaca}, yaitu memprediksi kata berikutnya dari citra.
Proses ini berjalan 200.000 langkah pada 64 GPU A100.

\textbf{Metode fusion.} \textbf{Tidak ada fusion eksplisit}, karena masukannya hanya citra.
Teks tidak diambil dari OCR, tetapi \enquote{dibaca} sendiri oleh model. Hubungan visual
dan bahasa terjadi lewat \textit{cross-attention} \textit{decoder} ke keluaran
\textit{encoder}.

\textbf{Hasil.}
\begin{itemize}
  \item RVL-CDIP: akurasi \textbf{95,30\%}, setara LayoutLMv2 (95,25\%), dengan waktu
  752~ms per dokumen berbanding 1.489~ms.
  \item CORD: F1 84,1.
  \item DocVQA: ANLS 67,5, masih di bawah LayoutLMv2 (78,1). Namun Donut lebih kuat pada
  tulisan tangan (ANLS 72,1).
\end{itemize}

Temuan paling relevan untuk topik kita ada di Gambar~9 paper. Akurasi dan kecepatan model
berbasis OCR (BERT dan LayoutLMv2) di CORD \textbf{sangat bervariasi tergantung mesin OCR
yang dipakai}. CLOVA OCR memberi akurasi terbaik.

\textbf{Hambatan dan \textit{future work}.}
\begin{itemize}
  \item Teks yang sangat kecil di halaman besar dapat terlewat karena batas resolusi
  masukan.
  \item Resolusi yang lebih tinggi menaikkan akurasi, tetapi biayanya membesar.
  \item Penulis menyebut perbaikan tujuan pralatih sebagai arah lanjutan.
\end{itemize}

\subsection{UDOP (CVPR 2023)}\label{sec:udop}

\textbf{Latar belakang.} Model sebelumnya memperlakukan citra dan teks dokumen sebagai dua
masukan terpisah. Tata letak hanya diwakili \textit{embedding} posisi yang dangkal, dan
setiap tugas butuh \textit{head} khusus. Padahal teks dokumen \enquote{tertanam} di citra
dengan korespondensi satu-satu \parencite{tang2023udop}.

\textbf{Dataset.}
\begin{itemize}
  \item \textbf{Pralatih:} IIT-CDIP (11 juta dokumen), ditambah 11 dataset berlabel
  (1,8 juta contoh), termasuk RVL-CDIP, PubLayNet, DocVQA, dan lainnya.
  \item \textbf{Evaluasi:} FUNSD, CORD, RVL-CDIP, dan tujuh dataset DUE-Benchmark.
\end{itemize}

\textbf{Teknologi OCR.}
\begin{itemize}
  \item IIT-CDIP dipakai beserta teks dan kotak hasil OCR yang tersedia, tanpa nama mesin.
  \item Untuk semua tugas evaluasi, penulis memakai \textbf{anotasi OCR asli yang disediakan
  dataset}.
\end{itemize}

\textbf{Pipeline.}
\begin{enumerate}
  \item Citra dipecah menjadi \textit{patch}.
  \item \textit{Embedding} setiap token teks \textbf{dijumlahkan} dengan fitur
  \textit{patch} citra tempat pusat kotak token itu berada (\textit{layout-induced
  vision-text embedding}).
  \item Kotak tata letak diubah menjadi token diskret, misalnya
  \texttt{<50><100><250><300>}.
  \item Semua tugas, termasuk klasifikasi, dirumuskan sebagai pembangkitan teks dengan
  prompt, misalnya \enquote{Document Classification on RVLCDIP}.
\end{enumerate}

\textbf{Model.} \textit{Vision-Text-Layout Transformer} dengan 794 juta parameter:
\begin{itemize}
  \item \textit{encoder} dan \textit{decoder} teks-tata letak berbasis T5-large;
  \item \textit{decoder} visual berbasis MAE-large.
\end{itemize}
Resolusi dinaikkan bertahap $224 \rightarrow 512 \rightarrow 1024$.

\textbf{Metode fusion.} Fusion paling awal. Teks dan citra dijumlahkan sudah di tingkat
\textit{embedding} masukan, berdasarkan letak spasialnya, lalu diproses satu
\textit{encoder}. Varian pembanding UDOP-Dual memakai dua \textit{encoder} terpisah untuk
teks dan citra (\textit{two-tower}).

\textbf{Hasil.}
\begin{itemize}
  \item RVL-CDIP: akurasi \textbf{96,00\%}; varian UDOP-Dual 96,22\%.
  \item FUNSD: F1 91,62. CORD: F1 97,58.
  \item Peringkat pertama di DUE-Benchmark per November 2022.
\end{itemize}

Menghapus input citra hanya menurunkan DocVQA dari 84,7 ke 84,4, dan InfographicsVQA dari
47,4 ke 45,0. Penulis menyimpulkan citra lebih penting pada data yang kaya visual daripada
pada data yang didominasi teks.

\textbf{Hambatan dan \textit{future work}.} Penulis menyebut dua hal:
\begin{itemize}
  \item kemampuan membangkitkan dokumen dapat disalahgunakan untuk pemalsuan;
  \item penerapan pada bahasa non-Latin memerlukan modifikasi.
\end{itemize}
Hambatan praktis yang tidak dibahas penulis: ukuran model (794 juta parameter) dan
pralatihnya sangat berat.

\subsection{DocFormerv2 (AAAI 2024)}\label{sec:docformerv2}

\textbf{Latar belakang.} Menderetkan teks OCR dari kiri ke kanan lalu memberikannya ke
model bahasa tidak optimal. Makna isi dokumen bergantung pada tata letaknya, dan sebagian
besar tugas butuh pemahaman \textbf{lokal}: hubungan sebuah teks dengan teks di
sekitarnya. Model sebelumnya kebanyakan hanya \textit{encoder}
\parencite{appalaraju2024docformerv2}.

\textbf{Dataset.}
\begin{itemize}
  \item \textbf{Pralatih:} Industrial Document Library (IDL), 64 juta halaman.
  \item \textbf{Evaluasi:} TabFact, WikiTableQuestions, DocVQA, InfographicsVQA, FUNSD, dan
  CORD, ditambah tiga dataset \textit{scene-text} VQA (OCR-VQA, TextVQA, ST-VQA).
  \item \textbf{RVL-CDIP tidak diuji}, sehingga tidak ada hasil klasifikasi dokumen.
\end{itemize}

\textbf{Teknologi OCR.} Teks dan kotak IDL diekstrak dengan OCR, tetapi \textbf{nama mesin
OCR-nya tidak disebutkan}. Untuk evaluasi dipakai teks OCR yang disediakan dataset.

\textbf{Pipeline.}
\begin{enumerate}
  \item Citra masuk lewat lapisan sederhana \texttt{linear(conv2x2(v))}, tanpa CNN atau
  detektor objek.
  \item Teks OCR ditokenisasi \textit{sentence-piece}.
  \item Fitur spasial (x, y, tinggi, lebar) dijumlahkan ke \textit{embedding} teks.
  \item Urutan teks dan citra disambung lalu masuk \textit{encoder}.
  \item \textit{Decoder} membangkitkan jawaban.
\end{enumerate}

\textbf{Model.} Transformer \textit{encoder-decoder}, varian \textit{small},
\textit{base}, dan \textit{large} (750 juta parameter). Pralatihnya dibuat asimetris:
\begin{itemize}
  \item dua tugas di \textit{encoder}: \textit{token-to-line} (jarak baris antar token) dan
  \textit{token-to-grid} (letak token di kisi halaman $4\times4$);
  \item satu tugas di \textit{decoder}: \textit{denoising language modeling}.
\end{itemize}

\textbf{Metode fusion.} Fusion awal dalam satu \textit{encoder}. Token visual dan token teks
disambung pada dimensi urutan, dan setiap modalitas diberi \textit{embedding} penanda
modalitas.

\textbf{Hasil.}
\begin{itemize}
  \item DocVQA: ANLS 87,84 dan InfographicsVQA 48,8. Keduanya setelah pelatihan tambahan
  dengan 850 ribu pasangan tanya-jawab dokumen.
  \item FUNSD: F1 88,89. CORD: F1 97,7.
  \item TabFact: 4,3 poin di atas model terbaik sebelumnya.
\end{itemize}

Dua ablasi relevan:
\begin{itemize}
  \item \textbf{Ketahanan terhadap galat OCR buatan di FUNSD.} Dengan 20\% karakter salah,
  kinerja DocFormerv2 hanya turun 1,68\%, sedangkan model \textit{encoder}-saja turun
  9,84\%.
  \item Memakai Swinv2 sebagai \textit{encoder} citra justru 4,3\% lebih buruk di DocVQA
  dibanding lapisan linear sederhana.
\end{itemize}

\textbf{Hambatan dan \textit{future work}.} Tidak disebutkan secara eksplisit. Dari analisis
kami: klasifikasi dokumen tidak diuji, mesin OCR pralatih tidak disebutkan, dan
pralatihnya memakai 64 juta halaman.

\subsection{DocLLM (ACL 2024)}\label{sec:docllm}

\textbf{Latar belakang.} Model bahasa besar (LLM) standar hanya menerima teks dan
mengandaikan tata letak sederhana. LLM multimodal untuk dokumen bergantung pada
\textit{encoder} citra yang mahal. DocLLM menambahkan informasi tata letak ke LLM
\textbf{tanpa \textit{encoder} citra} sama sekali, hanya dari kotak teks OCR
\parencite{wang2024docllm}.

\textbf{Dataset.}
\begin{itemize}
  \item \textbf{Pralatih:} IIT-CDIP dan DocBank, total 16,7 juta halaman (3,8 miliar
  token).
  \item \textbf{\textit{Instruction tuning}:} 16 dataset untuk empat tugas, yaitu
  tanya-jawab, inferensi bahasa alami, ekstraksi informasi, dan klasifikasi. Untuk
  klasifikasi dipakai RVL-CDIP (di-\textit{downsample}) dan BizDocs.
\end{itemize}

\textbf{Teknologi OCR.} Teks dan kotak berasal dari OCR yang disediakan setiap dataset;
mesinnya tidak disebutkan. Penulis secara eksplisit menyatakan bahwa \textbf{pilihan
mesin OCR masih terbuka untuk dieksplorasi}, dan perbandingan praktis antar-mesin OCR
dijadikan \textit{future work}. Mereka juga mencatat bahwa model keluarga LayoutLM yang
memakai Microsoft Azure OCR mengungguli Tesseract di papan peringkat DocVQA.

\textbf{Pipeline.}
\begin{enumerate}
  \item Setiap token teks OCR dipasangkan dengan kotaknya (kiri, atas, kanan, bawah).
  \item Kotak di-\textit{encode} menjadi vektor spasial terpisah dari vektor teks.
  \item Perhatian dihitung dengan mekanisme terurai (lihat bagian fusion).
  \item Tugas diberikan sebagai instruksi, misalnya \enquote{What type of document is this?
  Possible choices: \ldots}.
\end{enumerate}

\textbf{Model.} Dibangun di atas LLM yang sudah ada:
\begin{itemize}
  \item DocLLM-1B, berbasis Falcon-1B;
  \item DocLLM-7B, berbasis Llama2-7B.
\end{itemize}
Pralatihnya memakai tujuan \textit{block infilling}: blok teks di-\textit{mask} lalu
direkonstruksi dari konteks sebelum dan sesudahnya.

\textbf{Metode fusion.} Tidak ada citra; yang dipadukan adalah teks dan tata letak.
Caranya \textit{disentangled spatial attention}: skor perhatian dipecah menjadi teks--teks,
teks--spasial, spasial--teks, dan spasial--spasial, masing-masing dengan bobot $\lambda$.
Konfigurasi akhirnya hanya memakai teks--teks dan spasial--spasial.

\textbf{Hasil.}
\begin{itemize}
  \item RVL-CDIP: akurasi \textbf{91,8\%} (7B) dan 90,9\% (1B), jauh di atas GPT-4 +
  OCR \textit{zero-shot} (68,2\%).
  \item Unggul di 14 dari 16 dataset dibanding model sebanding.
\end{itemize}
Catatan: angka RVL-CDIP ini memakai format pilihan ganda dengan data latih yang
di-\textit{downsample}, sehingga tidak sebanding langsung dengan paper lain.

\textbf{Hambatan dan \textit{future work}.}
\begin{itemize}
  \item Pada dataset yang belum pernah dilihat, akurasi klasifikasinya rendah: 31,1\% pada
  BizDocs. Menurut penulis, ini karena model hanya dilatih dengan satu dataset
  klasifikasi.
  \item Penulis berencana menambahkan modalitas citra secara ringan.
  \item Perbandingan antar-mesin OCR disebut sebagai pekerjaan lanjutan.
\end{itemize}

\subsection{LiLT (ACL 2022)}\label{sec:lilt}

\textbf{Latar belakang.} Hampir semua model pemahaman dokumen terstruktur hanya untuk satu
bahasa, biasanya Inggris. Satu-satunya model multibahasa (LayoutXLM) butuh pengumpulan
data 53 bahasa yang mahal. LiLT memisahkan pengetahuan tata letak dari bahasa, sehingga
tata letak yang dipelajari dari dokumen Inggris dapat dipakai untuk bahasa lain
\parencite{wang2022lilt}.

\textbf{Dataset.}
\begin{itemize}
  \item \textbf{Pralatih:} IIT-CDIP (11 juta citra).
  \item \textbf{Evaluasi:} FUNSD, CORD, EPHOIE, RVL-CDIP, dan XFUND (delapan bahasa).
\end{itemize}

\textbf{Teknologi OCR.} Teks dan kotak IIT-CDIP diekstrak dengan \textbf{TextIn API}
(komersial). Setiap token memakai kotak dari baris teks tempatnya berada, karena tidak
semua mesin OCR memberikan kotak per token.

\textbf{Pipeline.}
\begin{enumerate}
  \item Teks OCR diurutkan dari kiri-atas ke kanan-bawah, lalu ditokenisasi.
  \item Kotak dinormalkan ke rentang 0--1000 menjadi urutan posisi 2D.
  \item Aliran teks dan aliran tata letak diproses dua Transformer paralel.
  \item Keluaran keduanya disambung untuk \textit{head} tugas.
\end{enumerate}
\textbf{Citra tidak dipakai.}

\textbf{Model.} Transformer dua aliran:
\begin{itemize}
  \item aliran teks diinisialisasi dari model bahasa siap pakai (RoBERTa, XLM-R, InfoXLM);
  \item aliran tata letak lebih kecil, hanya 6,1 juta parameter tambahan.
\end{itemize}
Pralatihnya memakai MVLM ditambah dua tugas baru: \textit{key point location} dan
\textit{cross-modal alignment identification}.

\textbf{Metode fusion.} \textit{Bi-directional attention complementation mechanism} (BiACM):
skor perhatian aliran teks dan aliran tata letak saling dijumlahkan di setiap lapis, lalu
fitur keduanya disambung di keluaran.

Ablasinya relevan untuk pilihan fusion. Menyambung fitur teks dan tata letak \textbf{hanya
di keluaran}, tanpa interaksi antar-aliran, \textbf{lebih buruk daripada aliran teks
saja}. \textit{Co-attention} yang lebih \enquote{dalam} juga menurunkan kinerja.

\textbf{Hasil.}
\begin{itemize}
  \item RVL-CDIP: akurasi \textbf{95,68\%} dengan RoBERTa Inggris.
  \item FUNSD: F1 88,41. CORD: F1 96,07.
  \item Unggul atas LayoutXLM di delapan bahasa XFUND, padahal hanya dipralatih dengan data
  Inggris.
\end{itemize}
Satu catatan: paper LayoutLMv3 menyebut LiLT memakai fitur citra ResNeXt101-FPN saat
\textit{fine-tuning} RVL-CDIP.

\textbf{Hambatan dan \textit{future work}.} Penulis akan:
\begin{itemize}
  \item melanjutkan pola transfer dari satu bahasa ke banyak bahasa;
  \item mengeksplorasi informasi visual yang tidak bergantung bahasa.
\end{itemize}

\subsection{Review Fusion untuk Klasifikasi Dokumen (Information Fusion 2026)}
\label{sec:mironczuk}

\textbf{Latar belakang.} Riset fusion untuk klasifikasi dokumen belum punya kerangka
bersama. Belum ada pula sintesis kuantitatif tentang seberapa besar sebenarnya keuntungan
fusion \parencite{mironczuk2026}.

\textbf{Dataset.} Bukan dataset eksperimen, melainkan 139 studi primer. Studi itu dipilih
dari 374 rekaman Scopus (pencarian 3 Januari 2025) dengan prosedur berbasis PRISMA.
Kriterianya:
\begin{itemize}
  \item berbahasa Inggris;
  \item membahas klasifikasi dokumen;
  \item memakai teknik fusion;
  \item telah ditinjau sejawat.
\end{itemize}

\textbf{Teknologi OCR.} Tidak dianalisis tersendiri. Penulis mengakui bahwa kualitas data,
termasuk galat OCR, adalah tantangan penting. Namun kualitas OCR \textbf{tidak dapat
dikontrol} dalam meta-analisis, karena studi primer jarang melaporkannya.

\textbf{Pipeline.}
\begin{enumerate}
  \item Identifikasi, penyaringan, dan seleksi studi.
  \item Sintesis kualitatif dan penyusunan taksonomi.
  \item Meta-analisis efek acak (\textit{random-effects}) atas selisih kinerja
  multimodal dibanding unimodal.
\end{enumerate}

\textbf{Model.} Kerangka formal representasi, pola, dan model, yang menghubungkan arsitektur
pembelajaran mesin dengan teori fusion klasik (Bayesian, Dempster--Shafer).

\textbf{Metode fusion.} Fusion dibagi menjadi tiga:
\begin{itemize}
  \item awal (\textit{early}), di tingkat representasi;
  \item akhir (\textit{late}), di tingkat keputusan;
  \item hibrida, gabungan keduanya.
\end{itemize}
Pada studi multimodal, pasangan teks--citra mendominasi (67,2\%).

\textbf{Hasil.}
\begin{itemize}
  \item Fusion multimodal menaikkan akurasi rata-rata \textbf{5,28 poin} ($p = 0{,}0016$).
  \item Efeknya terhadap F1 positif tetapi tidak signifikan.
  \item Hanya \textbf{11,8\%} studi multimodal yang memakai uji statistik.
  \item Keuntungan multimodal dapat mengecil bila sudah dipakai model bahasa pralatih yang
  kuat.
\end{itemize}

\textbf{Hambatan dan \textit{future work}.}
\begin{itemize}
  \item \textbf{Hambatan review ini:} hanya satu basis data (Scopus), dan kualitas data
  studi primer tidak terkontrol.
  \item \textbf{Hambatan di lapangan:} kelangkaan data, beban komputasi, penyelarasan
  lintas-modalitas, dan dokumen panjang.
  \item \textbf{Arah lanjutan:} integrasi dengan LLM, pembelajaran representasi
  kontrastif, penanganan data multimodal yang bising, dan validasi statistik yang lebih
  ketat.
\end{itemize}

\section{Perbandingan}\label{sec:perbandingan}

Tabel~\ref{tab:banding} merangkum aspek teknis keenam paper eksperimen. Tabel
\ref{tab:hambatan} merangkum hambatan dan \textit{future work}-nya.

\begin{table}[H]
  \centering
  \caption{Perbandingan teknis paper eksperimen}
  \label{tab:banding}
  \fontsize{9}{11}\selectfont
  \begin{tabularx}{\textwidth}{|l|>{\raggedright\arraybackslash}X|>{\raggedright\arraybackslash}X|>{\raggedright\arraybackslash}X|c|}
    \hline
    \textbf{Paper} & \textbf{OCR} & \textbf{Modalitas} & \textbf{Fusion} &
    \textbf{RVL-CDIP} \\ \hline
    LayoutLMv3 & Microsoft Read API & teks, tata letak, citra &
    awal, satu aliran & 95,93\% \\ \hline
    Donut & tidak ada (CLOVA OCR hanya untuk label pralatih) & citra &
    tidak ada; teks dibaca model & 95,30\% \\ \hline
    UDOP & bawaan dataset & teks, tata letak, citra &
    sangat awal, dijumlahkan per \textit{patch} & 96,00\% \\ \hline
    DocFormerv2 & tidak disebutkan & teks, tata letak, citra &
    awal, satu \textit{encoder} & tidak diuji \\ \hline
    DocLLM & bawaan dataset & teks, tata letak &
    perhatian terurai teks--spasial & 91,8\%$^\ast$ \\ \hline
    LiLT & TextIn API & teks, tata letak & dua aliran + BiACM & 95,68\% \\ \hline
  \end{tabularx}
  \sumber{Dirangkum dari paper masing-masing. $^\ast$Format pilihan ganda dengan data
  latih di-\textit{downsample}; tidak sebanding langsung.}
\end{table}

\begin{table}[H]
  \centering
  \caption{Hambatan dan \textit{future work} setiap paper}
  \label{tab:hambatan}
  \fontsize{9}{11}\selectfont
  \begin{tabularx}{\textwidth}{|l|>{\raggedright\arraybackslash}X|>{\raggedright\arraybackslash}X|}
    \hline
    \textbf{Paper} & \textbf{Hambatan} & \textbf{\textit{Future work} dari penulis} \\ \hline
    LayoutLMv3 & bergantung OCR komersial; pralatih sangat mahal &
    skala lebih besar; \textit{few-shot} dan \textit{zero-shot} \\ \hline
    Donut & teks kecil terlewat; resolusi tinggi mahal; DocVQA masih kalah dari model
    berbasis OCR & perbaikan tujuan pralatih \\ \hline
    UDOP & model sangat besar (794 juta parameter); risiko pemalsuan dokumen &
    adaptasi ke bahasa non-Latin \\ \hline
    DocFormerv2 & klasifikasi tidak diuji; OCR tidak disebutkan & tidak disebutkan \\ \hline
    DocLLM & klasifikasi lemah pada dataset baru & menambah citra secara ringan;
    \textbf{membandingkan mesin OCR} \\ \hline
    LiLT & citra tidak dipakai & transfer multibahasa; informasi visual \\ \hline
    Mirończuk & satu basis data; kualitas OCR tidak terkontrol & integrasi LLM;
    data bising; validasi statistik \\ \hline
  \end{tabularx}
\end{table}

\section{Sintesis dan Celah Penelitian}\label{sec:sintesis}

Dari ketujuh paper, ada empat pola yang terlihat.

\textbf{1. Mesin OCR diperlakukan sebagai konstanta.} Lima paper eksperimen bergantung pada
OCR, tetapi masing-masing memakai satu sumber saja:
\begin{itemize}
  \item Microsoft Read API (LayoutLMv3);
  \item TextIn API (LiLT);
  \item OCR bawaan dataset (UDOP, DocLLM);
  \item mesin yang tidak disebutkan (DocFormerv2).
\end{itemize}
Hanya Donut yang membandingkan beberapa mesin OCR, itu pun di lampiran dan untuk ekstraksi
informasi (CORD), bukan klasifikasi. Hasilnya kinerja sangat bergantung pada mesin OCR.
DocLLM secara eksplisit menyebut perbandingan mesin OCR sebagai \textit{future work}.
Mirończuk mencatat kualitas OCR tidak terkontrol di seluruh bidang.

\textbf{2. Sumbangan modalitas citra untuk klasifikasi kecil.}
\begin{itemize}
  \item Ablasi LayoutLMv3: teks dan tata letak saja sudah 95,33\% di RVL-CDIP; citra hanya
  menambah sekitar 0,2 poin.
  \item UDOP: menghapus citra hanya menurunkan DocVQA 0,3 poin.
  \item LiLT: tanpa citra sudah mencapai 95,68\%.
  \item Mirończuk: keuntungan multimodal mengecil bila sudah dipakai model bahasa yang
  kuat.
\end{itemize}
Tidak ada satu pun paper yang menguji apakah model gabungan \emph{benar-benar memakai}
setiap modalitas, misalnya lewat uji \textit{missing modality} di tingkat model.

\textbf{3. Cara fusion beragam.} Pilihannya meliputi fusion awal satu aliran (LayoutLMv3,
DocFormerv2), penjumlahan per \textit{patch} (UDOP), dua aliran dengan berbagi perhatian
(LiLT), perhatian terurai teks--spasial (DocLLM), sampai tanpa OCR sama sekali (Donut).
Ablasi LiLT menunjukkan bahwa penggabungan hanya di akhir tanpa interaksi bisa lebih buruk
daripada satu modalitas saja.

\textbf{4. Semuanya bergantung pada pralatih skala besar.} Data pralatihnya berkisar 11 juta
sampai 64 juta halaman, dengan puluhan GPU. Replikasi pralatih tidak realistis untuk
perangkat kelompok (GPU 4~GB). Jalur yang realistis adalah \textit{fine-tuning} model
\textit{base} yang dirilis publik pada data kecil seperti Tobacco3482, misalnya
LayoutLMv3-\textit{base} (133 juta parameter) atau LiLT.

\textbf{Celah yang terbuka.} Pola~1 dan~2 menunjuk ke pertanyaan yang belum dijawab paper
mana pun:
\begin{itemize}
  \item seberapa besar pilihan mesin OCR open-source (Tesseract, EasyOCR, PaddleOCR)
  memengaruhi model klasifikasi dokumen berbasis OCR;
  \item apakah model gabungan benar-benar memanfaatkan modalitas teks yang dihasilkan OCR
  tersebut.
\end{itemize}
Donut memberi bukti awal bahwa mesin OCR berpengaruh. DocLLM menyebut pertanyaan ini
sebagai pekerjaan lanjutan. Mirończuk menunjukkan bahwa celah ini terjadi di seluruh
bidang.
